\documentclass[11pt,a4paper]{ltjsarticle}
\usepackage{amsmath,amssymb}
\usepackage{graphicx}
\usepackage{booktabs}
\usepackage{array}
\usepackage{luatexja-fontspec}
\setmainjfont{HaranoAjiMincho-Regular}
\setsansjfont{HaranoAjiGothic-Medium}
\linespread{1.05}\selectfont

\begin{document}

\begin{center}
{\LARGE \textbf{視覚メディア処理I 第2回レポート}}
\end{center}
\vspace{1em}

\begin{flushright}
提出日: \today \\
学籍番号: 2611193 \\
氏名: 中井 平蔵
\end{flushright}

\section*{解答}

\subsection*{1. Zバッファ法とその特徴、長所・短所}

Zバッファ法とは、3次元物体を2次元画面に描画するときに、各画素に対して
「現在その画素に描かれている点の奥行き値」を記録しておき、手前にある面だけを
表示するための方法である。

具体的には、三角形ポリゴンをラスタライズするとき、各画素に対応する奥行き
$z$ を計算する。そして、その画素にすでに保存されている奥行き値 $z_{\mathrm{buf}}$
と比較する。以下の条件を満たす時、その画素の色と奥行き値を新しいものに更新する。
\[
z < z_{\mathrm{buf}}
\]

\subsubsection*{長所}
\begin{itemize}
  \item 実装が単純であり、各画素について奥行き値を比較するだけで隠面消去を行える。
  \item ポリゴンを奥行き順にソートする必要がないため、交差するポリゴンや複雑なシーンにも適用しやすい。
\end{itemize}

\subsubsection*{短所}
\begin{itemize}
  \item 各画素に奥行き値を保存するため、追加のメモリが必要となる。
  \item Zバッファは有限のビット数で奥行き値を表現するため、互いに非常に近い奥行きに位置する面では、前後関係の判定が不安定になり、Z-fighting が発生する場合がある。
  \item 透明物体の描画は、単純なZバッファだけでは正しく扱いにくい。透明面では奥の色も合成する必要があるため、追加処理が必要になる。
\end{itemize}

\subsection*{2. 三角形と線分の交差判定}

問題文より、3次元空間において、三角形 $F_1$ の３つの頂点は
\[
\mathbf{v}_1=(x_1,y_1,z_1)^\top,\quad
\mathbf{v}_2=(x_2,y_2,z_2)^\top,\quad
\mathbf{v}_3=(x_3,y_3,z_3)^\top
\]

である。また、線分 $L_1$ は原点から
\[
\mathbf{p}=(p_x,p_y,p_z)^\top
\]

までの線分である。この線分上の点は、パラメータ $t$ を用いて

\[
\mathbf{r}(t)=t\mathbf{p},\qquad 0\leq t\leq 1
\]

と書ける。

線分と三角形が交差するのは、線分上の点と三角形上の点が一致するときである。
三角形と同一平面上の点は、

\[
\mathbf{e}_1=\mathbf{v}_2-\mathbf{v}_1,\qquad
\mathbf{e}_2=\mathbf{v}_3-\mathbf{v}_1
\]

と重心座標 $u,v$ を用いて

\[
\mathbf{q}(u,v)=\mathbf{v}_1+u\mathbf{e}_1+v\mathbf{e}_2
\]

と表せる。この点が三角形の内部にある条件は、
\[
u\geq 0,\qquad v\geq 0,\qquad u+v\leq 1
\]

である。

したがって、線分と三角形が交差するかどうかは
\[
t\mathbf{p}=\mathbf{v}_1+u\mathbf{e}_1+v\mathbf{e}_2
\Rightarrow
\begin{bmatrix}\mathbf{p} & -\mathbf{e}_1 & -\mathbf{e}_2\end{bmatrix}\begin{bmatrix}t\\u\\v\end{bmatrix}=\mathbf{v}_1 \Rightarrow \begin{bmatrix}t\\u\\v\end{bmatrix}=\begin{bmatrix}\mathbf{p} & -\mathbf{e}_1 & -\mathbf{e}_2\end{bmatrix}^{-1}\mathbf{v}_1
\]

の解$t,u,v$が以下の条件を満たすかどうかを調べれば良い。

\[
0\leq t\leq 1,\qquad u\geq 0,\qquad v\geq 0,\qquad u+v\leq 1
\]

よって、判定手順は次のようになる。
\begin{enumerate}
  \item $\mathbf{e}_1=\mathbf{v}_2-\mathbf{v}_1$、$\mathbf{e}_2=\mathbf{v}_3-\mathbf{v}_1$ を計算する。
  \item 行列 $A=\begin{bmatrix}\mathbf{p}&-\mathbf{e}_1&-\mathbf{e}_2\end{bmatrix}$ を作り、$\det A$ を計算する。
  \item $\det A\approx 0$ の場合、線分が三角形の平面とほぼ平行であるため、\fbox{交差しない}と判定する。
  \item $\det A\neq 0$ の場合、連立方程式 $A(t,u,v)^\top=\mathbf{v}_1$ を解く。
  \item 得られた値が
  \[
  0\leq t\leq 1,\qquad u\geq 0,\qquad v\geq 0,\qquad u+v\leq 1
  \]
  をすべて満たすなら、線分 $L_1$ は三角形 $F_1$ と \fbox{交差する}。
  \item どれか1つでも満たさなければ、\fbox{交差しない}。
\end{enumerate}

\end{document}
